\section{Endpoint-access word support}\label{sec:endpoint}

The support restriction is clearest in a rectangular input--output setting.
Let
\[
A(\delta)=A_0+\sum_{r=1}^m\delta_rV_r,\qquad
M_k(\delta)=CA(\delta)^kB.
\]
For a multi-index \(\gamma\in\N^m\) with \(|\gamma|=j\), let
\(M_{\gamma,k}\) be the coefficient of \(\delta^\gamma\) in \(M_k\).
Define endpoint depths
\[
r_s(r)=\min\{q\ge0:V_rA_0^qB\ne0\},\qquad
r_t(r)=\min\{q\ge0:CA_0^qV_r\ne0\},
\]
with value \(\infty\) if the set is empty.  If a word realising \(\gamma\)
begins with label \(u\) at the source end and ends with label \(v\) at the
target end, call \((u,v)\) compatible.  Let
\[
\rho(\gamma)=\min_{(u,v)\ \mathrm{compatible}}
\bigl(r_s(u)+r_t(v)\bigr).
\]

\begin{theorem}[Endpoint-access bound]\label{thm:endpoint}
For \(j=|\gamma|\ge1\),
\begin{equation}\label{eq:endpointbound}
M_{\gamma,k}=0\quad\text{whenever}\quad k<j+\rho(\gamma).
\end{equation}
At \(k=j+\rho(\gamma)\), the coefficient is the sum of all compatible
endpoint-minimal words with zero internal gaps.  The bound is sharp if and only
if that boundary sum is non-zero.
\end{theorem}
\begin{proof}
Expand \(A(\delta)^k\) into labelled words.  A word contributing to
\(\delta^\gamma\) has \(j\) perturbation letters and \(k-j\) copies of \(A_0\).
If its first perturbation label at the source end is \(u\), at least
\(r_s(u)\) copies of \(A_0\) must lie to its right; if its last label at the
target end is \(v\), at least \(r_t(v)\) copies must lie to its left.  Internal
gaps are non-negative.  Therefore \(k-j\ge r_s(u)+r_t(v)\), and minimising gives
\cref{eq:endpointbound}.  Equality forces endpoint-minimal blocks and zero
internal gaps.  Summing all such words proves the boundary formula.  Distinct
words can cancel, so existence of words is not sharpness.
\end{proof}

The fully rectangular multi-index theorem above is \evidence{PROVED} by the
labelled-word argument.  The present Lean package formalises the locked
scalar-ray specialisation needed downstream, not this entire reusable
generality.

For the locked model \(B=e_0\), \(C=e_5^{\mathsf T}\), and
\(V_r=E_{rr}\) for \(r=1,\ldots,4\).  Every endpoint depth equals \(1\).

\begin{corollary}[Six-state support boundary]\label{cor:sixsupport}
For every \(\gamma\ne0\),
\[
[d^\gamma]M_k(d)=0\qquad(k<|\gamma|+2).
\]
At the minimal layer \(k=j+2\),
\begin{equation}\label{eq:minlayer}
\sum_{|\gamma|=j}[d^\gamma]M_{j+2}(d)d^\gamma
=d_1^j-d_2^j+d_3^j-d_4^j.
\end{equation}
Thus all mixed monomials vanish on the boundary.
\end{corollary}
\begin{proof}
The depth statement follows from the theorem.  On the boundary, the only word
is \(H_0D(d)^jH_0\), giving
\(e_5^{\mathsf T}H_0D(d)^jH_0e_0\).  Reading the four source and target edge
products in \cref{eq:H0} yields signs \(+,-,+,-\), proving
\cref{eq:minlayer}.
\end{proof}

\noindent\textbf{Formal status.}
For the locked channel, the complete finite word recurrence, the vanishing
below $k=j+2$, the boundary word and the signed power sum are
\evidence{KERNEL-CHECKED} as \path{M3A.rayCoeffMatrix},
\path{M3A.sixState_support_bound},
\path{M3A.rayMomentCoeff_boundary_matrix} and
\path{M3A.sixState_boundary_ray}.  These theorems sum actual word
coefficients, so boundary cancellation is already included.

The baseline \(j=0\) is deliberately excluded from the theorem: endpoint
depths are defined through a perturbation letter, and baseline darkness has its
own proof.  A finite census of 256 small labelled systems was used only to
stress-test common failure modes; it is \evidence{MEASURED} with \(n=256\), not
a proof of \cref{thm:endpoint}.
